\section{拘传}
\begin{enumerate}
  \item 唯一一种公检法都能决定并且执行。
  \item 主体。只能适用于\zx{非在押的嫌疑人/被告人}。公检 \yjt 嫌疑人。法院 \yjt 被告人。
  \item 持续时间不得超过 \zd{12 小时}，特别重大、复杂，可能拘留、速捕持续时间\zd{不得超过 24 小时}。
  \item \zd{保证休息时间}。\zd{2 次拘传间隔 \dydy 12 小时}。\al{连续讯问 24 小时 \cuo}。
  \item 在他所在市县的地点内进行、2 人以上。
  \item \bs{拘传证/拘传票没有例外}。
        \al{对现场发现的犯罪嫌疑人，经出示工作证件\xhx{可以口头}拘传，并在笔录中注明 \cuo}。
  \item \zx{立即讯问}。
  \item \zd{vs 传唤}。\zx{主体扩大到当事人、传票、可以口头}。其他基本一样。\al{传唤证人 \cuo}。
        \begin{enumerate}
          \item 是否为强制措施。
          \item 是否具有强制力。
          \item 是否需要法律文书。
        \end{enumerate}
\end{enumerate}
